%=============================================================================
%  Active Inference for Port Selection in Fluid Antenna Systems Under Partial CSI
%
%  This is a FAITHFUL LaTeX port of paper/working_draft.docx (generated by
%  paper/make_working_draft.py). The prose is Kian's, carried over verbatim.
%  The equations were already LaTeX strings in that script -- they were
%  rasterized to PNG only because python-docx cannot typeset math -- so they
%  are reproduced here as real, numbered, cross-referenced equations.
%
%  SCOPE NOTE: this draft predates partial sensing (m < M), the AR(p) temporal
%  model, and hybrid transmit. The upgraded versions of Sections II-III live in
%  sections_v2.tex, which is NOT input here. Merge them section by section.
%
%  BUILD:  ./build.sh
%=============================================================================
\documentclass[journal]{IEEEtran}

\usepackage{amsmath,amssymb,amsfonts}
\usepackage{algorithm}
\usepackage{algorithmic}
\usepackage{graphicx}
\usepackage{tikz}
\usetikzlibrary{arrows.meta,positioning,calc,fit,backgrounds,shapes.geometric}
\usepackage{textcomp}
\usepackage{xcolor}
\usepackage[colorlinks=true,allcolors=blue]{hyperref}

% ---- editorial markers: delete before submission -------------------------
\newcommand{\todo}[1]{\textcolor{red}{\textbf{[TODO: #1]}}}
\newcommand{\note}[1]{\textcolor{blue}{\textit{[#1]}}}

% ---- notation shorthands -------------------------------------------------
\newcommand{\bh}{\mathbf{h}}
\newcommand{\bH}{\mathbf{H}}
\newcommand{\bR}{\mathbf{R}}
\newcommand{\bI}{\mathbf{I}}
\newcommand{\bW}{\mathbf{W}}
\newcommand{\bP}{\mathbf{P}}
\newcommand{\bF}{\mathbf{F}}
\newcommand{\by}{\mathbf{y}}
\newcommand{\bn}{\mathbf{n}}
\newcommand{\bw}{\mathbf{w}}
\newcommand{\bz}{\mathbf{z}}
\newcommand{\bB}{\mathbf{B}}
\newcommand{\bmu}{\boldsymbol{\mu}}
\newcommand{\bSig}{\boldsymbol{\Sigma}}
\newcommand{\CN}{\mathcal{CN}}
\newcommand{\E}{\mathbb{E}}
\newcommand{\bbC}{\mathbb{C}}
\newcommand{\bbR}{\mathbb{R}}
\DeclareMathOperator{\tr}{tr}
\DeclareMathOperator*{\argmin}{arg\,min}
\DeclareMathOperator*{\argmax}{arg\,max}

\begin{document}

\title{Active Inference for Joint Port Selection and\\
       Pilot Allocation in Fluid Antenna Systems}

\author{Kian~Fotovat$^{1}$ and Kamran~Fotovat$^{2}$%
  \thanks{$^{1}$Dept.\ of Electrical and Computer Engineering, University of
  Tehran, Tehran, Iran.}%
  \thanks{$^{2}$Iran University of Science and Technology, Tehran, Iran.}%
}

\maketitle

%==============================================================================
% NOTE: numbers below are from results_tm/paper1_config.txt (MC=6). They are
% being re-measured at MC=20 by sim_version3/make_fig2_data.py; update here and
% in Section VI when that lands.
\begin{abstract}
Fluid antenna systems (FAS) obtain spatial degrees of freedom by activating a
few radiating ports from a dense grid of candidates, but deciding which ports
to activate presupposes channel knowledge at ports that are never measured.
Prior port-selection work assumes full or near-full channel state information
(CSI) and treats the pilot overhead as a fixed cost. We instead formulate port
selection and pilot allocation \emph{jointly} as active inference. A
spatio-temporal generative model --- sub-wavelength spatial correlation
together with an order-$p$ autoregressive temporal prior --- maintains a
Gaussian belief over all $N$ port channels, and the transmitter chooses both
the ports it serves and the subset of them it pilots by minimizing an expected
free energy that trades achievable rate against information gain and
port-switching cost. Selection is greedy and submodular, at $\mathcal{O}(NM)$
per slot. On a $441$-port grid with a hybrid front end at
$n_{\mathrm{RF}}=2K$, the agent attains 91\% of a full-CSI genie's sum rate
while measuring only 2.3\% of the candidate ports; pilots can then be withdrawn
from ports that are still served, one fifth of them for 2\% of the sum rate and
two fifths for 8\%. No training is required.
\end{abstract}

\begin{IEEEkeywords}
Fluid antenna systems, port selection, pilot allocation, active inference,
expected free energy, partial CSI, Kalman filtering.
\end{IEEEkeywords}

%==============================================================================
\section{Introduction}
\label{sec:intro}

Fluid antenna systems (FAS), also called movable- or reconfigurable-position
antenna systems, replace fixed radiating elements with elements whose position
can be adjusted among a dense set of candidate positions, or ports, within a
compact aperture \cite{wong2021fas}. By moving to a favourable point of the
spatial fading pattern, an FAS harvests diversity that a fixed array of the
same size cannot, using only a handful of radio-frequency (RF) chains, since at
any instant only the activated ports are connected to the front end. The
pivotal control problem is port selection: deciding slot by slot which subset
of the candidates to activate.

A port can be scored only if its channel is known, so conventional selection
presupposes channel state information (CSI) at every candidate. That is
expensive. The aperture is sampled at sub-wavelength spacing, so the ports are
packed far more densely than a half-wavelength array and the number of
coefficients to estimate grows rapidly with aperture. Worse, only activated
ports can be observed at all. The bottleneck is therefore not the selection
arithmetic but channel acquisition under a sensing budget: the system must
decide which ports are worth measuring --- and, we add, whether a port that is
served needs to be measured at all.

Two lines of work approach this and leave the gap open. Learning-based
selection --- deep reinforcement learning over ports \cite{drl_switching} and
online bandits \cite{online_portsel} --- learns policies directly, but treats
the full channel as an observed, cost-free input and requires extensive
training. Channel estimation for FAS \cite{sbar,fas_lmmse} exploits spatial
correlation to reconstruct the channel from few pilots, but solves estimation
in isolation: it recovers the channel without deciding what to sense or what to
serve, and never weighs a measurement against its cost. What is missing is a
formulation in which sensing and serving are chosen jointly under an explicit
budget.

Active inference supplies that perspective. Rooted in the free-energy principle
\cite{parr2022active}, it casts a decision-maker as an agent holding a
generative model of its environment and acting to minimize an expected free
energy (EFE), which decomposes into a pragmatic value rewarding goal attainment
and an epistemic value rewarding informative actions. The fit to FAS is close:
the hidden state is the channel at unmeasured ports, the generative model is
its spatial and temporal correlation, the pragmatic value is throughput, and
the epistemic value is what a pilot reveals about the ports left unseen. The
exploit--explore tension that defines the problem is thus expressed by the EFE
rather than engineered in by hand.

\begin{enumerate}
\item We formulate port selection \emph{and pilot allocation} jointly as active
inference under partial CSI. The piloted set is a strict subset of the
activated set, making the pilot budget a decision variable rather than a fixed
cost --- to our knowledge the first such treatment for FAS.

\item We give a concrete agent: a Gaussian belief over all $N$ ports from a
spatial prior and an order-$p$ autoregressive temporal model, with ports and
pilots chosen by greedy EFE minimization at $\mathcal{O}(NM)$ per slot, and an
exact reduced-rank belief that makes a $21\times21$ grid tractable.

\item We show that the agent reaches 91\% of a full-CSI genie's sum rate while
measuring $2.3\%$ of the ports, and degrades gracefully as pilots are further
withdrawn from ports it continues to serve.
\end{enumerate}

%==============================================================================
\section{Fluid Antenna System Model}
\label{sec:model}

\begin{figure}[t]
  \centering
  % =============================================================================
%  Fig. 1 -- per-slot protocol timeline.
%
%  The point of this figure: the ACTIVATED set S_t and the PILOTED set Q_t are
%  NOT the same set. Q_t is a strict subset. Without this picture a reviewer
%  will read m and M as the same quantity, which erases the contribution.
%
%  COLOUR. Palette is Okabe-Ito (colour-blind safe) and semantically tied:
%  BLUE always means "activated / S_t", ORANGE always means "piloted / Q_t".
%  The same two colours carry from the port grid into the timeline boxes, so a
%  reader can follow one idea by colour alone. Everything else stays neutral
%  grey so the two that matter stand out. Luminances are far enough apart that
%  the figure still reads if printed greyscale.
%
%  Geometry: laid out for a 3.5in (8.9cm) IEEE column; all coordinates are cm.
%  PREVIEW:  ./build_fig.sh fig_timeline
%  NOTE:     preview goes through pdflatex, never latex+dvipng -- dvipng
%            silently discards TikZ's PostScript specials. See build_fig.sh.
% =============================================================================
\definecolor{cAct}{HTML}{1F6FB2}     % activated ports, S_t
\definecolor{cActBg}{HTML}{E7F0F8}
\definecolor{cPil}{HTML}{D55E00}     % piloted ports, Q_t
\definecolor{cPilBg}{HTML}{FCEFE4}
\definecolor{cCand}{HTML}{C3CAD3}    % candidate ports
\definecolor{cNeut}{HTML}{5E6A75}    % neutral stage borders
\definecolor{cNeutBg}{HTML}{F2F4F7}
\definecolor{cArrow}{HTML}{4A5560}

\begin{tikzpicture}[
  x=1cm, y=1cm,
  font=\scriptsize,
  >={Stealth[length=1.5mm]},
  stage/.style={
    draw=cNeut, fill=cNeutBg, rounded corners=1.2pt, align=center,
    inner sep=1.5pt, minimum height=6.6mm, minimum width=15.5mm,
    line width=0.45pt},
  actbox/.style={stage, draw=cAct, fill=cActBg, line width=0.6pt},
  pilbox/.style={stage, draw=cPil, fill=cPilBg, line width=0.6pt},
  sub/.style={font=\fontsize{5.6}{6.6}\selectfont, align=center, inner sep=1pt},
  key/.style={font=\fontsize{5.8}{6.6}\selectfont, inner sep=1pt, anchor=west},
]

% ------------------------------------------------------------ port grid
% 7x7 schematic of the N candidate ports. Five are activated (S_t); three of
% those five are also piloted (Q_t).
\begin{scope}[shift={(0,2.55)}]
  \foreach \c in {0,...,6}
    \foreach \r in {0,...,6}
      \fill[cCand] (0.24*\c, 0.24*\r) circle (0.030);

  % activated ports S_t   (|S_t| = M)
  \foreach \p in {(0.24,1.20), (0.96,1.44), (0.48,0.48), (1.20,0.72), (0.72,0)}
    \fill[cAct] \p circle (0.062);

  % piloted subset Q_t    (|Q_t| = m) -- ringed
  \foreach \p in {(0.24,1.20), (1.20,0.72), (0.72,0)}
    \draw[cPil, line width=0.7pt] \p circle (0.125);

  % an activated-but-unpiloted port, for the callout to point at
  \coordinate (unpiloted) at (0.48,0.48);
\end{scope}

% ------------------------------------------------------------ legend
\fill[cCand] (2.05,3.87) circle (0.030);
\node[key] at (2.20,3.87) {one of $N$ candidate ports};

\fill[cAct] (2.05,3.54) circle (0.062);
\node[key] at (2.20,3.54)
  {\textcolor{cAct}{\textbf{activated}}, $\mathcal{S}_t$, $|\mathcal{S}_t|=M$
   --- carries data};

\fill[cAct] (2.05,3.21) circle (0.062);
\draw[cPil, line width=0.7pt] (2.05,3.21) circle (0.125);
\node[key] at (2.20,3.21)
  {\textcolor{cPil}{\textbf{piloted}}, $\mathcal{Q}_t\subseteq\mathcal{S}_t$,
   $|\mathcal{Q}_t|=m\le M$};

% ------------------------------------------------------------ callout
\node[font=\fontsize{5.8}{7}\selectfont, anchor=north west,
      text width=6.15cm, align=left, inner sep=1pt] (callout) at (2.20,2.60)
  {the $M-m$ activated-but-unpiloted ports are \emph{served without ever being
   measured}, inferred through $\mathbf{R}$ and the AR($p$) prior};

\draw[->, line width=0.4pt, cPil]
  (callout.west) to[out=180, in=-20] ($(unpiloted)+(0.11,-0.02)$);

% ------------------------------------------------------------ timeline
\node[stage]  (predict) at (0.875,1.10) {predict\\[-1pt]belief};
\node[actbox] (select)  at (2.575,1.10) {select $\mathcal{S}_t$};
\node[pilbox] (sense)   at (4.275,1.10) {pilot $\mathcal{Q}_t$};
\node[stage]  (update)  at (5.975,1.10) {update\\[-1pt]belief};
\node[stage]  (tx)      at (7.675,1.10) {precode\\[-1pt]\& transmit};

\draw[->, cArrow] (predict) -- (select);
\draw[->, cArrow] (select)  -- (sense);
\draw[->, cArrow] (sense)   -- (update);
\draw[->, cArrow] (update)  -- (tx);

\node[sub, below=1.0mm of predict] {AR($p$)\\[-1pt]dynamics};
\node[sub, below=1.0mm of select, text=cAct]
  {greedy EFE\\[-1pt]$\mathcal{O}(NM)$};
\node[sub, below=1.0mm of sense, text=cPil]
  {\emph{uplink}\\[-1pt]$m$ pilots};
\node[sub, below=1.0mm of update]  {Kalman\\[-1pt]all $N$ ports};
\node[sub, below=1.0mm of tx]      {\emph{downlink}\\[-1pt]$M$ ports};

% belief carried into the next slot
\draw[->, dashed, line width=0.4pt, cNeut]
  (tx.north) -- (7.675,1.78)
  -- node[sub, above, pos=0.5, text=cNeut] {belief carried to slot $t+1$}
  (0.875,1.78) -- (predict.north);

\end{tikzpicture}

  \caption{Per-slot operation. The base station activates $M$ of $N$ candidate
  ports, but spends pilots on only $m\le M$ of them; the remaining $M-m$
  activated ports carry data while never being measured, their channels
  supplied by the belief alone. Both the activated set $\mathcal{S}_t$ and the
  piloted subset $\mathcal{Q}_t$ are chosen by minimizing the same expected
  free energy, which is what turns the pilot budget from a fixed cost into a
  design variable.}
  \label{fig:timeline}
\end{figure}

We consider the downlink of a single-cell multiuser system in which a base
station (BS) equipped with a fluid antenna serves $K$ single-antenna users.
This section specifies the array geometry, the spatio-temporal channel, the
partial observation model that is central to this work, the transmission and
rate model, and the switching-aware objective.

\subsection{Array Geometry and Activation}
The BS fluid antenna comprises $N=N_x\times N_y$ candidate ports arranged on a
uniform planar grid over a compact aperture with sub-half-wavelength spacing.
The BS has $M$ radio-frequency (RF) chains, $M\le N$, so in each slot $t$ it
activates a subset $\mathcal{S}(t)\subseteq\{1,\dots,N\}$ of
$|\mathcal{S}(t)|=M$ ports; only activated ports radiate and connect to the
front end. Serving $K$ streams requires $M\ge K$, hence $N\ge M\ge K$. We write
$\bP_{\mathcal{S}}\in\{0,1\}^{M\times N}$ for the selection matrix whose rows
are the standard basis vectors indexed by $\mathcal{S}$, so that
$\bP_{\mathcal{S}}\mathbf{x}$ extracts the entries of a vector $\mathbf{x}$ on
the activated ports.

\subsection{Spatio-Temporal Channel Model}
Let $\bh_k(t)\in\bbC^{N}$ denote the channel between the $N$ candidate ports
and user $k$ in slot $t$. Because the ports are densely spaced, their
coefficients are spatially correlated. Following the classical Jakes model for
rich scattering, the correlation between ports $i$ and $j$ depends only on
their separation $d_{ij}$,
\begin{equation}
  R_{ij} = J_0\!\left(\frac{2\pi d_{ij}}{\lambda}\right),
  \label{eq:R}
\end{equation}
where $J_0(\cdot)$ is the zeroth-order Bessel function of the first kind and
$\lambda$ the wavelength. Collecting these entries into $\bR\in\bbC^{N\times
N}$, the small-scale fading is spatially correlated Rayleigh,
$\bh_k\sim\CN(\mathbf{0},\beta_k\bR)$, with $\beta_k$ the large-scale gain of
user $k$. The channel ages between slots according to a first-order
Gauss--Markov (autoregressive) process,
\begin{equation}
  \bh_k(t) = \rho\,\bh_k(t-1)+\sqrt{1-\rho^2}\,\mathbf{e}_k(t),
  \quad
  \mathbf{e}_k(t)\sim\CN(\mathbf{0},\beta_k\bR),
  \label{eq:ar1}
\end{equation}
where $\mathbf{e}_k(t)$ is independent of $\bh_k(t-1)$, and the temporal
correlation $\rho=J_0(2\pi f_D T_s)$ is set by the maximum Doppler frequency
$f_D$ and the slot duration $T_s$; $\rho\to1$ for slow fading and decreases
with mobility. The process is stationary with marginal
$\bh_k(t)\sim\CN(\mathbf{0},\beta_k\bR)$.

\subsection{Observation Model: Partial CSI}
A defining constraint of our setting is that the BS cannot observe the full
channel: information about a port is obtained only when that port is activated.
During the pilot phase of slot $t$ the users send orthogonal pilots and the BS
estimates the channel on the $M$ activated ports, while the remaining $N-M$
ports are unobserved. The measurement is
\begin{equation}
  \by_k(t) = \bP_{\mathcal{S}(t)}\,\bh_k(t)+\bn_k(t),
  \quad
  \bn_k(t)\sim\CN(\mathbf{0},\sigma_e^2\bI_M),
  \label{eq:obs}
\end{equation}
where $\bP_{\mathcal{S}(t)}$ restricts $\bh_k(t)$ to the activated ports and
$\bn_k(t)$ is the pilot/estimation noise with per-port variance $\sigma_e^2$
set by the pilot SNR. Equivalently, the BS observes the length-$M$ vector
$\bh_{k,\mathcal{S}}(t)=\bP_{\mathcal{S}(t)}\bh_k(t)$ in noise and must infer
the $N-M$ hidden ports from spatial and temporal correlation. This partial
observability distinguishes our formulation from prior port-selection work that
assumes full-channel knowledge, and it is what makes the choice of which ports
to measure consequential.

\subsection{Downlink Transmission and Achievable Rate}
Over the activated ports the BS transmits with a linear precoder
$\bW(t)=[\bw_1(t),\dots,\bw_K(t)]\in\bbC^{M\times K}$, one column per user,
under a total power budget $\|\bW(t)\|_F^2\le P$. The transmitted signal is
$\mathbf{x}(t)=\sum_k\bw_k(t)s_k(t)$ with data symbols
$s_k\sim\CN(0,1)$. User $k$ receives
\begin{equation}
  r_k = \bh_{k,\mathcal{S}}^{\mathrm{H}}\bw_k s_k
      + \sum_{j\neq k}\bh_{k,\mathcal{S}}^{\mathrm{H}}\bw_j s_j + z_k,
  \label{eq:rx}
\end{equation}
where $\bh_{k,\mathcal{S}}=\bP_{\mathcal{S}}\bh_k$ is the channel on the
activated ports and $z_k\sim\CN(0,\sigma^2)$ is receiver noise. Treating
inter-user interference as noise, the SINR and achievable rate of user $k$ are
\begin{equation}
  \mathrm{SINR}_k =
  \frac{\big|\bh_{k,\mathcal{S}}^{\mathrm{H}}\bw_k\big|^2}
       {\sum_{j\neq k}\big|\bh_{k,\mathcal{S}}^{\mathrm{H}}\bw_j\big|^2+\sigma^2},
  \label{eq:sinr}
\end{equation}
\begin{equation}
  R_k = \log_2\!\left(1+\mathrm{SINR}_k\right).
  \label{eq:rate}
\end{equation}
The precoder is obtained by regularized MMSE (transmit-Wiener) filtering from
the available channel estimate
$\hat{\bH}=[\hat{\bh}_{1,\mathcal{S}},\dots,\hat{\bh}_{K,\mathcal{S}}]
\in\bbC^{M\times K}$,
\begin{equation}
  \bW = \xi\,\hat{\bH}
    \left(\hat{\bH}^{\mathrm{H}}\hat{\bH}+\frac{K\sigma^2}{P}\bI_K\right)^{-1},
  \label{eq:mmse}
\end{equation}
where $\xi$ scales $\bW$ to meet the power budget. Under full CSI, $\hat{\bH}$
is the true activated channel; in our partial-CSI setting $\hat{\bH}$ and its
error statistics are supplied by the belief of Section~\ref{sec:aif}, yielding
a robust MMSE precoder that additionally accounts for the residual channel
uncertainty.

\subsection{Switching Cost and Problem Formulation}
Reconfiguring the fluid antenna between slots is not free: physically moving or
re-tuning elements incurs delay and energy. We model this by the number of
ports that change between consecutive slots,
\begin{equation}
  C_t = \big|\mathcal{S}_t\,\triangle\,\mathcal{S}_{t-1}\big|,
  \qquad
  E_{\mathrm{sw}}(t) = e_{\mathrm{sw}}\,C_t,
  \label{eq:switch}
\end{equation}
where $\triangle$ is the symmetric set difference and $e_{\mathrm{sw}}$ the
per-port switching energy. The design goal is to choose the activation sequence
that maximizes long-term throughput net of switching cost,
\begin{equation}
  \max_{\{\mathcal{S}_t\}}\;
  \sum_{t=1}^{T}\left(\sum_{k=1}^{K}R_k(t)
    -\eta_{\mathrm{sw}}E_{\mathrm{sw}}(t)\right)
  \;\;\text{s.t.}\;\;|\mathcal{S}_t|=M,
  \label{eq:problem}
\end{equation}
where $\eta_{\mathrm{sw}}\ge0$ trades throughput against reconfiguration cost.
Equation~\eqref{eq:problem} is the objective against which every method in this
paper is evaluated. Its difficulty is that $R_k(t)$ depends on the channel over
ports that, under partial CSI, are largely unobserved --- the problem we
address next through active inference.

%==============================================================================
\section{Active Inference for Port Selection}
\label{sec:aif}

We now cast port selection as active inference. The agent holds a generative
model of the channel, maintains a Bayesian belief over all ports (perception),
and each slot chooses the port set that minimizes an expected free energy
(action), which trades expected rate against information gain and switching
cost within one objective.

\subsection{Generative Model}
The latent state is the full channel $\{\bh_k(t)\}$, of which only the
activated ports are observed. The agent's generative model factorizes over
users and time as
\begin{equation}
  \begin{split}
    p\big(\{\bh_k,\by_k\}\big) = \prod_{k}p\big(\bh_k(0)\big)\prod_{t}\;
      & p\big(\bh_k(t)\,|\,\bh_k(t-1)\big)\\[-2pt]
      & \times\;p\big(\by_k(t)\,|\,\bh_k(t),\mathcal{S}_t\big),
  \end{split}
  \label{eq:gen}
\end{equation}
with three ingredients already fixed by the physics of
Section~\ref{sec:model}: the prior $p(\bh_k(0))=\CN(\mathbf{0},\beta_k\bR)$;
the transition $p(\bh_k(t)|\bh_k(t-1))$ given by the AR(1) dynamics
\eqref{eq:ar1}; and the likelihood $p(\by_k(t)|\bh_k(t),\mathcal{S}_t)$ given
by the partial observation \eqref{eq:obs}. Preferences complete the model.
Unlike textbook active inference, in which preferred outcomes are encoded in a
hand-tuned vector, here the preference is the communication utility itself ---
the achievable rate net of switching in \eqref{eq:problem}. There is thus no
arbitrary preference to set: the agent's goal is the system objective.

\subsection{Perception: Complex Kalman Belief}
Because the model is linear-Gaussian, the exact posterior over each user's
channel is Gaussian, $q(\bh_k)=\CN(\bmu_k,\bSig_k)$, maintained by a per-user
complex Kalman filter with two steps per slot. The predict step propagates the
belief through the AR(1) dynamics, encoding channel aging,
\begin{equation}
  \bmu_k \leftarrow \rho\,\bmu_k,
  \qquad
  \bSig_k \leftarrow \rho^2\bSig_k+(1-\rho^2)\beta_k\bR .
  \label{eq:predict}
\end{equation}
When the ports in $\mathcal{S}_t$ are activated and $\by_k$ observed, the
measurement update corrects the belief. With Kalman gain
\begin{equation}
  \mathbf{K}_k = \bSig_k\bP_{\mathcal{S}}^{\mathrm{H}}
    \left(\bP_{\mathcal{S}}\bSig_k\bP_{\mathcal{S}}^{\mathrm{H}}
      +\sigma_e^2\bI_M\right)^{-1},
  \label{eq:gain}
\end{equation}
the posterior mean and covariance become
\begin{equation}
  \bmu_k \leftarrow \bmu_k
    + \mathbf{K}_k\big(\by_k-\bP_{\mathcal{S}}\bmu_k\big),
  \label{eq:mean}
\end{equation}
\begin{equation}
  \bSig_k \leftarrow
    \big(\bI-\mathbf{K}_k\bP_{\mathcal{S}}\big)\bSig_k
    \big(\bI-\mathbf{K}_k\bP_{\mathcal{S}}\big)^{\mathrm{H}}
    + \sigma_e^2\mathbf{K}_k\mathbf{K}_k^{\mathrm{H}} .
  \label{eq:joseph}
\end{equation}
where the Joseph form \eqref{eq:joseph} keeps $\bSig_k$ symmetric positive
semidefinite. Only the $M\times M$ innovation covariance is inverted, and it is
regularized by $\sigma_e^2\bI$; consequently the rank-deficient spatial prior
$\bR$ needs no artificial jitter. Crucially, observing a port also shrinks the
uncertainty of its correlated neighbours through the off-diagonals of $\bR$ and
$\bSig_k$, so a few measurements inform many ports. The filter is calibrated:
its predicted covariance matches the realized estimation error.

\subsection{Action: Expected Free Energy}
Each slot the agent scores a candidate port set $\mathcal{S}$ by its expected
free energy, written here as a value to be minimized,
\begin{equation}
  G(\mathcal{S}) = -\,\mathrm{Prag}(\mathcal{S})
    -\kappa\,\mathrm{Epis}(\mathcal{S})
    +\eta_{\mathrm{sw}}E_{\mathrm{sw}}(\mathcal{S}),
  \label{eq:efe}
\end{equation}
a pragmatic (goal) term, an epistemic (information) term weighted by an
exploration weight $\kappa\ge0$, and the switching cost of \eqref{eq:switch}.
The pragmatic value is the expected sum rate the belief predicts for
$\mathcal{S}$ under the robust MMSE precoder; accounting for residual CSI error
through the belief covariance gives the imperfect-CSI lower bound
\begin{equation}
  \mathrm{Prag}(\mathcal{S}) = \sum_{k}
    \log_2\!\left(1+
      \frac{\big|\bmu_{k,\mathcal{S}}^{\mathrm{H}}\bw_k\big|^2}
           {\mathcal{I}_k}\right),
  \label{eq:prag}
\end{equation}
\begin{equation}
  \mathcal{I}_k = \sum_{j\neq k}
      \big|\bmu_{k,\mathcal{S}}^{\mathrm{H}}\bw_j\big|^2
    + \sum_{j}\bw_j^{\mathrm{H}}\bSig_{k,\mathcal{S}}\bw_j
    + \sigma^2,
  \label{eq:leak}
\end{equation}
where $\bmu_{k,\mathcal{S}}=\bP_{\mathcal{S}}\bmu_k$,
$\bSig_{k,\mathcal{S}}=\bP_{\mathcal{S}}\bSig_k\bP_{\mathcal{S}}^{\mathrm{H}}$,
and the error term $\sum_j\bw_j^{\mathrm{H}}\bSig_{k,\mathcal{S}}\bw_j$ makes
the agent conservative when uncertain; the robust MMSE precoder augments
\eqref{eq:mmse} with $\bSig_{k,\mathcal{S}}$. The epistemic value is the mutual
information between the channel and the pilots that $\mathcal{S}$ would yield,
i.e., the expected reduction in belief entropy,
\begin{equation}
  \mathrm{Epis}(\mathcal{S}) = \sum_{k}
    \log_2\det\!\left(\bI_M
      + \frac{\bP_{\mathcal{S}}\bSig_k\bP_{\mathcal{S}}^{\mathrm{H}}}
             {\sigma_e^2}\right),
  \label{eq:epis}
\end{equation}
which is monotone and submodular in $\mathcal{S}$; through the correlation in
$\bSig_k$, activating one port also gathers information about others. The two
terms embody the exploit--explore tension natively: the pragmatic value favours
ports that serve well now, the epistemic value favours ports that most reduce
uncertainty for later, and $\kappa$ sets the balance --- no separate
exploration heuristic is bolted on.

\subsection{Observe-Then-Precode Protocol}
The order of operations within a slot matters. We adopt an observe-then-precode
protocol: (i) predict the belief; (ii) select $\mathcal{S}_t$ by minimizing $G$
on the predicted belief; (iii) activate $\mathcal{S}_t$, obtain the pilots
$\by_k$, and update the belief; and (iv) form the robust MMSE precoder from the
updated belief and transmit. Because the served ports are re-measured before
precoding, their posterior error is about $\sigma_e^2$ rather than the one-step
aging floor $(1-\rho^2)\beta_k$ that a predict-then-act ordering incurs. As
shown in Section~\ref{sec:results}, this makes the served rate nearly
independent of Doppler.

\subsection{Greedy Submodular Selection}
Selecting the best $M$-of-$N$ set exactly is combinatorial. We instead build
$\mathcal{S}_t$ greedily: starting from the empty set, we repeatedly add the
port whose inclusion most decreases $G$,
\begin{equation}
  p^{\star} = \argmax_{p\notin\mathcal{S}}\;
    \Big[\,G(\mathcal{S})-G\big(\mathcal{S}\cup\{p\}\big)\,\Big],
  \label{eq:greedy}
\end{equation}
until $|\mathcal{S}|=M$. Since the epistemic term is monotone submodular and
the switching term modular, greedy selection attains the standard $(1-1/e)$
near-optimality guarantee on the information-driven part, at
$\mathcal{O}(NM)$ score evaluations per slot --- about 20\,ms in our setting
versus seconds for exhaustive search, a roughly $500\times$ reduction.
Algorithm~\ref{alg:aif} summarizes the complete per-slot agent, with the greedy
loop (lines 3--8) embedded in the observe-then-precode ordering of
Section~\ref{sec:aif}-D.

\begin{algorithm}[t]
\caption{Active-Inference Port Selection (slot $t$)}
\label{alg:aif}
\begin{algorithmic}[1]
\REQUIRE predicted belief $\{\bmu_k,\bSig_k\}$, previous set
         $\mathcal{S}_{t-1}$, budget $M$, weight $\kappa$
\ENSURE  activated set $\mathcal{S}_t$ and precoder $\bW$
\STATE predict belief: $\bmu_k\leftarrow\rho\bmu_k$,
       $\bSig_k\leftarrow\rho^2\bSig_k+(1-\rho^2)\beta_k\bR$ $\;\forall k$
\STATE $\mathcal{S}\leftarrow\emptyset$
\FOR{$m=1$ \TO $M$}
  \FOR{each candidate port $p\notin\mathcal{S}$}
    \STATE $\Delta(p)\leftarrow G(\mathcal{S})-G(\mathcal{S}\cup\{p\})$
           \hfill $\triangleright$ via \eqref{eq:efe}--\eqref{eq:epis}
  \ENDFOR
  \STATE $p^{\star}\leftarrow\argmax_p\Delta(p)$;\;
         $\mathcal{S}\leftarrow\mathcal{S}\cup\{p^{\star}\}$
\ENDFOR
\STATE $\mathcal{S}_t\leftarrow\mathcal{S}$
\STATE activate $\mathcal{S}_t$, observe pilots $\by_k$, update belief
       \hfill $\triangleright$ via \eqref{eq:gain}--\eqref{eq:joseph}
\STATE $\bW\leftarrow$ robust MMSE precoder from updated
       $\{\bmu_k,\bSig_k\}$
\STATE transmit; incur switching cost
       $\eta_{\mathrm{sw}}\cdot|\mathcal{S}_t\,\triangle\,\mathcal{S}_{t-1}|$
\end{algorithmic}
\end{algorithm}

\subsection{Online Learning of the Spatial Correlation}
The agent assumes a Jakes $\bR$, but real propagation may differ. When $\bR$ is
unknown or non-Jakes it can be learned online from the accumulated pilots:
averaging the outer products of co-observed pilots and removing the noise floor
gives
\begin{equation}
  \hat{R}_{ij} = \frac{\sum_t \mathbf{1}\{i,j\in\mathcal{S}_t\}\,
                       y_i(t)\,y_j^{*}(t)}
                      {\sum_t \mathbf{1}\{i,j\in\mathcal{S}_t\}}
                 - \sigma_e^2\,\delta_{ij},
  \label{eq:learnR}
\end{equation}
which is then normalized to unit diagonal (cancelling the per-user power
$\beta_k$) and projected onto the positive-semidefinite cone. Substituting
$\hat{\bR}$ for $\bR$ restores performance under model mismatch
(Section~\ref{sec:results}); in the matched case the belief is already
Bayes-optimal, so learning can only help --- consistent with the method's
zero-training nature.

%==============================================================================
\section{Numerical Results}
\label{sec:results}
\todo{Not yet written in the working draft. Figures exist in figures/.}

%==============================================================================
\section{Conclusion}
\label{sec:conc}
\todo{Not yet written in the working draft.}

%==============================================================================
\bibliographystyle{IEEEtran}
\bibliography{refs}

\end{document}
